% ===================================================================
% PHẦN NÀY DO: Trần Hoàng Phúc phụ trách
% Nội dung: Triển khai RTSP/RTP cơ bản (4 điểm)
% ===================================================================

\section{Triển khai giao thức RTSP ở Client và đóng gói RTP ở Server (4 điểm)}

\textbf{Phần này chỉ bao gồm triển khai CƠ BẢN của giao thức RTSP và RTP, KHÔNG bao gồm các tính năng mở rộng như frame fragmentation cho HD video.}

\subsection{RTSP Protocol - Client Side}

\subsubsection{RTSP State Machine \& Constants}
Client quản lý trạng thái phiên RTSP thông qua state machine với 3 trạng thái chính:

\begin{lstlisting}[language=Python, caption={RTSP State Machine}]
class Client:
    # RTSP States
    INIT = 0      # Trang thai khoi tao
    READY = 1     # San sang phat video
    PLAYING = 2   # Dang phat video
    state = INIT
    
    # RTSP Request Types
    SETUP = 0
    PLAY = 1
    PAUSE = 2
    TEARDOWN = 3
\end{lstlisting}

\textbf{Luồng chuyển trạng thái:}
\begin{itemize}
    \item \texttt{INIT → READY}: Sau khi nhận SETUP response thành công
    \item \texttt{READY → PLAYING}: Sau khi nhận PLAY response thành công
    \item \texttt{PLAYING → READY}: Sau khi nhận PAUSE response thành công
    \item \texttt{READY/PLAYING → INIT}: Sau khi nhận TEARDOWN response thành công
\end{itemize}

\subsubsection{Gửi RTSP Requests}

\paragraph{Hàm sendRtspRequest()}
Hàm này xử lý việc tạo và gửi các RTSP request đến server qua kết nối TCP.

\begin{lstlisting}[language=Python, caption={SETUP Request}]
def sendRtspRequest(self, requestCode):
    """Gui RTSP request den server"""
    request = ""
    
    if requestCode == self.SETUP and self.state == self.INIT:
        # Bat dau thread nhan RTSP responses
        threading.Thread(target=self.recvRtspReply, daemon=True).start()
        
        self.rtspSeq += 1
        # Format: SETUP <filename> RTSP/1.0
        request = f"SETUP {self.fileName} RTSP/1.0\n"
        request += f"CSeq: {self.rtspSeq}\n"
        request += f"Transport: RTP/UDP; client_port={self.rtpPort}"
        self.requestSent = self.SETUP
\end{lstlisting}

\textbf{Chi tiết SETUP Request:}
\begin{itemize}
    \item Khởi tạo thread \texttt{recvRtspReply} để nhận response bất đồng bộ
    \item Tăng \texttt{rtspSeq} để đánh số thứ tự request
    \item Header \texttt{Transport} chỉ định cổng UDP mà client sẽ nhận RTP packets
    \item Lưu loại request đã gửi vào \texttt{requestSent}
\end{itemize}

\begin{lstlisting}[language=Python, caption={PLAY Request}]
elif requestCode == self.PLAY and self.state == self.READY:
    self.rtspSeq += 1
    # Format: PLAY <filename> RTSP/1.0
    request = f"PLAY {self.fileName} RTSP/1.0\n"
    request += f"CSeq: {self.rtspSeq}\n"
    request += f"Session: {self.sessionId}"
    self.requestSent = self.PLAY
\end{lstlisting}

\textbf{Chi tiết PLAY Request:}
\begin{itemize}
    \item Chỉ được gửi khi ở trạng thái \texttt{READY}
    \item Bao gồm \texttt{Session ID} nhận được từ SETUP response
    \item Không có header \texttt{Transport} (đã thiết lập ở SETUP)
\end{itemize}

\begin{lstlisting}[language=Python, caption={PAUSE Request}]
elif requestCode == self.PAUSE and self.state == self.PLAYING:
    self.rtspSeq += 1
    # Format: PAUSE <filename> RTSP/1.0
    request = f"PAUSE {self.fileName} RTSP/1.0\n"
    request += f"CSeq: {self.rtspSeq}\n"
    request += f"Session: {self.sessionId}"
    self.requestSent = self.PAUSE
\end{lstlisting}

\textbf{Chi tiết PAUSE Request:}
\begin{itemize}
    \item Chỉ được gửi khi đang ở trạng thái \texttt{PLAYING}
    \item Tạm dừng streaming nhưng giữ nguyên phiên kết nối
    \item Có thể tiếp tục bằng PLAY request
\end{itemize}

\begin{lstlisting}[language=Python, caption={TEARDOWN Request \& Sending}]
elif requestCode == self.TEARDOWN and not self.state == self.INIT:
    self.rtspSeq += 1
    # Format: TEARDOWN <filename> RTSP/1.0
    request = f"TEARDOWN {self.fileName} RTSP/1.0\n"
    request += f"CSeq: {self.rtspSeq}\n"
    request += f"Session: {self.sessionId}"
    self.requestSent = self.TEARDOWN
else:
    return

try:
    # Gui qua RTSP socket (TCP)
    self.rtspSocket.send(request.encode("utf-8"))
    print(f"Sent RTSP request:\n{request}")
except Exception as e:
    print(f"Send error: {e}")
\end{lstlisting}

\textbf{Chi tiết TEARDOWN Request:}
\begin{itemize}
    \item Kết thúc phiên streaming và giải phóng tài nguyên
    \item Đóng tất cả socket connections
    \item Đưa client về trạng thái \texttt{INIT}
\end{itemize}

\subsubsection{Nhận RTSP Responses}

\paragraph{Hàm recvRtspReply()}
Thread này chạy liên tục để nhận và xử lý RTSP responses từ server.

\begin{lstlisting}[language=Python, caption={Receive RTSP Reply Thread}]
def recvRtspReply(self):
    """Thread nhan RTSP responses tu server"""
    while True:
        try:
            reply = self.rtspSocket.recv(1024)
            if reply:
                self.parseRtspReply(reply.decode("utf-8"))
            if self.requestSent == self.TEARDOWN:
                self.rtspSocket.shutdown(socket.SHUT_RDWR)
                self.rtspSocket.close()
                break
        except:
            break
\end{lstlisting}

\textbf{Chức năng:}
\begin{itemize}
    \item Nhận response từ RTSP socket (TCP)
    \item Decode và parse response
    \item Đóng socket sau khi nhận TEARDOWN response
\end{itemize}

\paragraph{Hàm parseRtspReply()}
Parse RTSP response và cập nhật trạng thái client.

\begin{lstlisting}[language=Python, caption={Parse RTSP Response}]
def parseRtspReply(self, data):
    """Parse RTSP response va cap nhat state"""
    try:
        lines = data.split('\n')
        seqNum = int(lines[1].split(' ')[1])
        
        # Verify sequence number
        if seqNum == self.rtspSeq:
            session = int(lines[2].split(' ')[1])
            
            # Luu session ID lan dau
            if self.sessionId == 0:
                self.sessionId = session
            
            # Verify session ID
            if self.sessionId == session:
                # Kiem tra status code
                if int(lines[0].split(' ')[1]) == 200:
                    if self.requestSent == self.SETUP:
                        # SETUP OK -> READY state
                        self.state = self.READY
                        self.openRtpPort()
                        # Bat dau nhan RTP packets
                        threading.Thread(target=self.listenRtp, 
                                       daemon=True).start()
                        
                    elif self.requestSent == self.PLAY:
                        # PLAY OK -> PLAYING state
                        self.state = self.PLAYING
                        
                    elif self.requestSent == self.PAUSE:
                        # PAUSE OK -> READY state
                        self.state = self.READY
                        self.playEvent.set()
                        
                    elif self.requestSent == self.TEARDOWN:
                        # TEARDOWN OK -> INIT state
                        self.state = self.INIT
                        self.teardownAcked = 1
    except Exception as e:
        print(f"Parse error: {e}")
\end{lstlisting}

\textbf{Xử lý response:}
\begin{itemize}
    \item Xác thực \texttt{CSeq} và \texttt{Session ID}
    \item Kiểm tra status code (200 = OK)
    \item Cập nhật state machine dựa trên loại request
    \item Khởi tạo RTP listener sau SETUP thành công
\end{itemize}

\subsubsection{Kết nối RTSP \& RTP}

\begin{lstlisting}[language=Python, caption={Connect to Server}]
def connectToServer(self):
    """Tao RTSP socket (TCP) va ket noi den server"""
    self.rtspSocket = socket.socket(socket.AF_INET, socket.SOCK_STREAM)
    try:
        self.rtspSocket.connect((self.serverAddr, self.serverPort))
        print(f"RTSP connection established")
    except Exception as e:
        print(f"Connection failed: {e}")

def openRtpPort(self):
    """Mo UDP socket de nhan RTP packets"""
    try:
        self.rtpSocket = socket.socket(socket.AF_INET, socket.SOCK_DGRAM)
        self.rtpSocket.settimeout(0.5)
        self.rtpSocket.bind(('', self.rtpPort))
        print(f"RTP port {self.rtpPort} opened")
    except Exception as e:
        print(f"Unable to bind RTP port: {e}")
\end{lstlisting}

\textbf{Hai loại kết nối:}
\begin{itemize}
    \item \textbf{RTSP Socket (TCP):} Dùng cho control messages
    \item \textbf{RTP Socket (UDP):} Dùng cho streaming data với timeout 0.5s
\end{itemize}

\subsection{RTP Encoding - Server Side}

\subsubsection{RTP Header Structure}

RTP packet được thiết kế theo chuẩn RFC 3550 với header 12 bytes.

\paragraph{Header Layout (RFC 3550)}

\begin{verbatim}
 0                   1                   2                   3
 0 1 2 3 4 5 6 7 8 9 0 1 2 3 4 5 6 7 8 9 0 1 2 3 4 5 6 7 8 9 0 1
+-+-+-+-+-+-+-+-+-+-+-+-+-+-+-+-+-+-+-+-+-+-+-+-+-+-+-+-+-+-+-+-+
|V=2|P|X|  CC   |M|     PT      |       sequence number         |
+-+-+-+-+-+-+-+-+-+-+-+-+-+-+-+-+-+-+-+-+-+-+-+-+-+-+-+-+-+-+-+-+
|                           timestamp                           |
+-+-+-+-+-+-+-+-+-+-+-+-+-+-+-+-+-+-+-+-+-+-+-+-+-+-+-+-+-+-+-+-+
|           synchronization source (SSRC) identifier            |
+-+-+-+-+-+-+-+-+-+-+-+-+-+-+-+-+-+-+-+-+-+-+-+-+-+-+-+-+-+-+-+-+
|                            payload                            |
+-+-+-+-+-+-+-+-+-+-+-+-+-+-+-+-+-+-+-+-+-+-+-+-+-+-+-+-+-+-+-+-+
\end{verbatim}

\textbf{Header Size:} 12 bytes (96 bits)

\subsubsection{Field Definitions}

\begin{table}[H]
\centering
\begin{tabular}{|l|c|l|c|}
\hline
\textbf{Field} & \textbf{Bits} & \textbf{Description} & \textbf{Value} \\
\hline
V (Version) & 2 & RTP version & 2 \\
P (Padding) & 1 & Padding flag & 0 \\
X (Extension) & 1 & Extension flag & 0 \\
CC (CSRC Count) & 4 & CSRC count & 0 \\
M (Marker) & 1 & Marker bit & 0 \\
PT (Payload Type) & 7 & Payload type & 26 (MJPEG) \\
Sequence Number & 16 & Packet sequence & Tăng dần \\
Timestamp & 32 & Sampling instant & Unix time \\
SSRC & 32 & Sync source & 0 \\
\hline
\end{tabular}
\caption{RTP Header Fields}
\end{table}

\textbf{Lưu ý:} Trong triển khai cơ bản, Marker bit = 0 (mỗi frame được gửi trong 1 packet). Marker bit sẽ được sử dụng cho frame fragmentation trong phần HD Video Streaming.

\subsubsection{Decode RTP Packets}

\begin{lstlisting}[language=Python, caption={RTP Decode Methods}]
def decode(self, byteStream):
    """Decode RTP packet"""
    self.header = bytearray(byteStream[:HEADER_SIZE])
    self.payload = byteStream[HEADER_SIZE:]

def version(self):
    return int(self.header[0] >> 6)

def seqNum(self):
    seqNum = self.header[2] << 8 | self.header[3]
    return int(seqNum)

def timestamp(self):
    timestamp = (self.header[4] << 24 | self.header[5] << 16 | 
                self.header[6] << 8 | self.header[7])
    return int(timestamp)

def payloadType(self):
    pt = self.header[1] & 127
    return int(pt)

def marker(self):
    """Return Marker bit"""
    return int(self.header[1] >> 7)

def getPayload(self):
    return self.payload
    
def getPacket(self):
    """Return complete RTP packet (header + payload)"""
    return self.header + self.payload
\end{lstlisting}

\subsubsection{Tạo RTP Packets (Cơ bản)}

\begin{lstlisting}[language=Python, caption={Make RTP Packet - Basic Version}]
def makeRtp(self, payload, frameNbr):
    """
    Tao RTP packet co ban (1 frame = 1 packet)
    
    Args:
        payload: Frame data
        frameNbr: Sequence number
    
    Returns:
        Complete RTP packet (header + payload)
    """
    version = 2      # RTP version 2
    padding = 0      # No padding
    extension = 0    # No extension
    cc = 0          # No CSRC
    marker = 0      # Marker bit = 0 trong triển khai cơ bản
    pt = 26         # Payload Type: MJPEG (RFC 2435)
    seqnum = frameNbr
    ssrc = 0        # Synchronization source identifier
    
    rtpPacket = RtpPacket()
    rtpPacket.encode(version, padding, extension, cc, 
                     seqnum, marker, pt, ssrc, payload)
    return rtpPacket.getPacket()
\end{lstlisting}

\textbf{Tham số RTP packet cơ bản:}
\begin{itemize}
    \item \textbf{Version = 2:} Chuẩn RTP version 2
    \item \textbf{PT = 26:} MJPEG payload type theo RFC 2435
    \item \textbf{Marker bit = 0:} Trong triển khai cơ bản, mỗi frame là 1 packet
    \item \textbf{Sequence number:} Tăng dần theo số frame
\end{itemize}

\subsubsection{Gửi RTP Packets (Cơ bản)}

\begin{lstlisting}[language=Python, caption={Send RTP Stream - Basic Version}]
def sendRtp(self):
    """Main streaming loop - gui RTP packets (co ban)"""
    while True:
        frame_start = time.time()
        
        # Check for pause/stop
        if self.clientInfo['event'].isSet():
            print("Stopping RTP stream")
            break
        
        # Doc frame tu video file
        data = self.clientInfo['videoStream'].nextFrame()
        
        if not data:
            print("End of video")
            break
        
        frameNumber = self.clientInfo['videoStream'].frameNbr()
        
        try:
            address = self.clientInfo['rtspSocket'][1][0]
            port = int(self.clientInfo['rtpPort'])
            
            # Tao va gui 1 RTP packet cho 1 frame
            packet = self.makeRtp(data, frameNumber)
            self.clientInfo['rtpSocket'].sendto(packet, (address, port))
            
            self.frames_sent += 1
            self.bytes_sent += len(data)
            
        except Exception as e:
            print(f"Send error: {e}")
            break
        
        # Frame rate control (25 FPS)
        frame_elapsed = time.time() - frame_start
        time.sleep(max(0, 0.04 - frame_elapsed))
\end{lstlisting}

\textbf{Streaming loop cơ bản:}
\begin{itemize}
    \item Đọc frame từ video file
    \item Tạo 1 RTP packet cho toàn bộ frame (không chia nhỏ)
    \item Gửi qua UDP
    \item Frame rate control: 25 FPS (0.04s/frame)
    \item Dừng khi nhận pause/stop event
\end{itemize}

\subsubsection{Xử lý RTSP Requests}

\begin{lstlisting}[language=Python, caption={Process RTSP Request - SETUP}]
def processRtspRequest(self, data):
    """Xu ly RTSP requests tu client"""
    try:
        request = data.split('\n')
        line1 = request[0].split(' ')
        requestType = line1[0]
        filename = line1[1]
        seq = request[1].split(' ')
        
        if requestType == self.SETUP:
            if self.state == self.INIT:
                print("Processing SETUP")
                try:
                    # Mo video file
                    self.clientInfo['videoStream'] = VideoStream(filename)
                    self.state = self.READY
                    
                    # Tao UDP socket cho RTP
                    self.clientInfo["rtpSocket"] = socket.socket(
                        socket.AF_INET, socket.SOCK_DGRAM
                    )
                    
                except IOError:
                    self.replyRtsp(self.FILE_NOT_FOUND_404, seq[1])
                    return
                
                # Generate random session ID
                self.clientInfo['session'] = randint(100000, 999999)
                
                # Gui RTSP reply
                self.replyRtsp(self.OK_200, seq[1])
                
                # Lay RTP port tu request
                self.clientInfo['rtpPort'] = request[2].split('=')[1]
\end{lstlisting}

\textbf{SETUP processing:}
\begin{itemize}
    \item Mở video file
    \item Tạo UDP socket cho RTP streaming
    \item Generate random Session ID
    \item Gửi 200 OK response với Session ID
\end{itemize}

\begin{lstlisting}[language=Python, caption={Process RTSP Request - PLAY/PAUSE/TEARDOWN}]
elif requestType == self.PLAY:
    if self.state == self.READY:
        print("processing PLAY\n")
        self.state = self.PLAYING
        self.start_time = time.time()
        
        self.replyRtsp(self.OK_200, seq[1])
        
        # Bat dau streaming
        self.clientInfo['event'] = threading.Event()
        self.clientInfo['worker'] = threading.Thread(
            target=self.sendRtp, 
            daemon=True
        )
        self.clientInfo['worker'].start()

elif requestType == self.PAUSE:
    if self.state == self.PLAYING:
        print("processing PAUSE\n")
        self.state = self.READY
        self.clientInfo['event'].set()
        self.replyRtsp(self.OK_200, seq[1])

elif requestType == self.TEARDOWN:
    print("processing TEARDOWN\n")
    self.clientInfo['event'].set()
    self.replyRtsp(self.OK_200, seq[1])
    
    # Cleanup
    if 'videoStream' in self.clientInfo:
        self.clientInfo['videoStream'].close()
    if 'rtpSocket' in self.clientInfo:
        self.clientInfo['rtpSocket'].close()
\end{lstlisting}

\begin{lstlisting}[language=Python, caption={Reply RTSP}]
def replyRtsp(self, code, seq):
    """Gui RTSP reply"""
    if code == self.OK_200:
        reply = (f'RTSP/1.0 200 OK\n'
                f'CSeq: {seq}\n'
                f'Session: {self.clientInfo["session"]}')
        connSocket = self.clientInfo['rtspSocket'][0]
        connSocket.send(reply.encode())
        print(f"Sent RTSP reply:\n{reply}")
    elif code == self.FILE_NOT_FOUND_404:
        print("404 FILE NOT FOUND")
    elif code == self.CON_ERR_500:
        print("500 CONNECTION ERROR")
\end{lstlisting}

\textbf{RTSP responses:}
\begin{itemize}
    \item \textbf{200 OK:} Request thành công
    \item \textbf{404 Not Found:} Video file không tồn tại
    \item \textbf{500 Error:} Lỗi kết nối
\end{itemize}

\subsection{Tóm tắt triển khai RTSP/RTP cơ bản}

\subsubsection{RTSP Client}
\begin{itemize}
    \item State machine: INIT → READY → PLAYING
    \item 4 methods: SETUP, PLAY, PAUSE, TEARDOWN
    \item TCP connection cho RTSP control messages
    \item UDP socket cho RTP data reception
    \item Async thread để nhận responses và RTP packets
\end{itemize}

\subsubsection{RTP Server}
\begin{itemize}
    \item RTP header 12 bytes theo RFC 3550
    \item \textbf{1 frame = 1 RTP packet} (không chia nhỏ frame)
    \item Marker bit = 0 trong triển khai cơ bản
    \item Sequence number tăng dần theo frame number
    \item UDP transmission với frame rate 25 FPS
    \item Session management với random Session ID
\end{itemize}

\textbf{Lưu ý quan trọng:} Phần triển khai cơ bản này chỉ hỗ trợ video với frame size nhỏ (có thể gửi trong 1 UDP packet). Để hỗ trợ HD video với frame size lớn, cần triển khai frame fragmentation (xem Section 4).

\clearpage
